\documentclass[10pt]{article}
\usepackage{graphicx}
\usepackage{booktabs}
\usepackage{amsmath}
\usepackage{url}
\setlength{\parindent}{0pt}
\setlength{\parskip}{6pt}
\setlength{\columnsep}{12pt}
\setlength{\textwidth}{16.0cm}
\setlength{\oddsidemargin}{0.5cm}
\setlength{\evensidemargin}{0.5cm}
\setlength{\topmargin}{-0.5cm}
\setlength{\headheight}{12pt}

\title{Metabolic State Clustering Analysis: Choudhury et al. 2022 (3CA:20773)}
\author{}
\date{}

\begin{document}

\sloppy

\maketitle

\section{Introduction}

This analysis investigates whether transcriptional metabolic states can form distinct clusters in the Choudhury et al. 2022 single-cell RNA-seq dataset (3CA:20773). The question addressed is: \textit{Are there distinct clusters of transcriptional metabolic states? In other words, can we use these metabolic genes only to cluster or separate groups in the single-cell RNA-seq dataset?}

\section{Methods}

\subsection{Dataset}

The analysis uses the Choudhury et al. 2022 brain single-cell RNA-seq dataset from the Weizmann Institute 3CA catalog (3CA:20773). The dataset includes:

\begin{itemize}
    \item Expression matrix: 33,538 genes $\times$ 58,843 cells
    \item Technology: 10x Genomics
    \item Disease: Meningioma
    \item Samples: 10
\end{itemize}

\subsection{Normalization and Feature Selection}

Expression data was normalized using per-cell library size scaling to 10,000 counts, followed by log1p transformation:

\begin{equation}
    \label{eq:normalization}
    x_{ij}^{\text{norm}} = \log_2\left(1 + \frac{x_{ij}}{L_j}\right) \times 10,000
\end{equation}

where $x_{ij}$ is the raw UMI count for gene $i$ in cell $j$, $L_j$ is the library size for cell $j$, and the final scaling factor of 10,000 normalizes expression values. Equation~\ref{eq:normalization} defines the normalization formula.

\subsection{Feature Scaling}

Features were scaled using Scanpy's scale function with zero\_center=False and max\_value=10.

\subsection{Clustering}

Leiden clustering was performed on the metabolic gene set (1,984 genes after prevalence filtering from 2,197 Reactome genes). Resolution selection was performed across candidates [0.4, 0.8, 1.2] with stability filtering (mean pairwise ARI $\ge$ 0.75) and tie-breaking by mean pairwise ARI.

\subsection{Baselines and Controls}

\begin{itemize}
    \item HVG baseline: Highest variance of library-normalized log1p expression among prevalence-passing non-mitochondrial symbols
    \item Random controls: 19 size-matched samples without replacement from prevalence-passing non-metabolic, non-mitochondrial symbols
    \item KMeans control: 15 clusters with seed 0
\end{itemize}

\section{Results}

\subsection{Global Clustering Analysis}

\begin{figure}[h!]
    \centering
    \includegraphics{../figures/metabolic_umap.pdf}
    \caption{UMAP embedding of 58,843 cells clustered using Leiden algorithm on 1,984 metabolic genes. Cluster sizes range from 1,002 to 8,985 cells across 15 clusters.}
    \label{fig:metabolic_umap}
\end{figure}

Figure~\ref{fig:metabolic_umap} shows the UMAP embedding of cells using the metabolic gene set. The clustering identified 15 distinct groups with varying sizes, ranging from 1,002 to 8,985 cells.

\subsection{Clustering Diagnostics}

\begin{figure}[h!]
    \centering
    \includegraphics{../figures/clustering_diagnostics.pdf}
    \caption{Clustering diagnostics including silhouette scores, Calinski-Harabasz index, Davies-Bouldin index, and pairwise ARI/NMI metrics across seeds and resolutions.}
    \label{fig:clustering_diagnostics}
\end{figure}

Figure~\ref{fig:clustering_diagnostics} shows the clustering diagnostics. The selected resolution (0.4) maximized median silhouette among candidates passing the mean pairwise ARI stability filter; mean pairwise ARI was the tie-breaker.

\subsection{Clustering Metrics}

\begin{table}[h!]
    \centering
    \caption{Clustering metrics for the metabolic gene set analysis}
    \label{tab:clustering_metrics}
    \begin{tabular}{lr}
        \toprule
        Metric & Value \\
        \midrule
        Leiden silhouette & 0.229 \\
        Calinski-Harabasz & 7,177.0 \\
        Davies-Bouldin & 1.513 \\
        Mean pairwise ARI & 0.980 \\
        Min pairwise ARI & 0.973 \\
        Mean pairwise NMI & 0.975 \\
        Min pairwise NMI & 0.968 \\
        Median silhouette & 0.221 \\
        \bottomrule
    \end{tabular}
\end{table}

Table~\ref{tab:clustering_metrics} summarizes the clustering metrics. The silhouette score of 0.229 indicates moderate clustering quality.

\subsection{Random Control Comparison}

The observed random-control exceedance count was 0 of 19; the add-one rank fraction is descriptive, not an inferential p-value. Clustering performance was evaluated across seeds [0, 17, 42, 73, 101] with resolution candidates [0.4, 0.8, 1.2].

\subsection{Confounder Analysis}

The following confounders were checked: patient, sample, source, cell\_type, cell\_subtype, cell\_cycle\_phase, complexity, computed\_total\_counts, computed\_detected\_genes, and computed\_mito\_percent.

Categorical confounder NMI values:
\begin{itemize}
    \item patient: 0.544
    \item sample: 0.567
    \item source: 0.197
    \item cell\_type: 0.507
    \item cell\_subtype: 0.658
    \item cell\_cycle\_phase: 0.195
\end{itemize}

\subsection{Within-Replicate Sensitivity}

The analysis assessed within-replicate sensitivity across 6 patients. The weakest replicate was patient 3 with 3,287 cells. Global vs within-replicate ARI ranged from 0.106 to 0.813.

\section{Discussion}

\subsection{Question 1: Distinct Metabolic States}

The analysis found candidate partitions of observed groups but \textbf{distinct metabolic states remain inconclusive}. The clustering identified 15 groups, but the random control comparison (0 of 19 random controls exceeding metabolic clustering) suggests these partitions may not represent biologically distinct metabolic states.

\subsection{Question 2: Cell-Type Associations}

The associations between clusters and cell types/subtypes are \textbf{descriptive and non-causal}. The NMI values indicate moderate associations, but these should not be interpreted as causal relationships.

\subsection{Question 3: CD8 T Cell Subtype Analysis}

The CD8 T cell subtype analysis also found candidate partitions but \textbf{distinct metabolic states remain inconclusive}.

\section{Limitations}

\begin{itemize}
    \item Cell-level tests do not create biological replicates
    \item The analysis is exploratory and does not test discrete biological states against a continuum
    \item No automatic threshold proves independence, multimodality, mechanism, causality, or statistical significance
    \item Cell counts are not independent biological replicates
\end{itemize}

\section{References}

\begin{thebibliography}{1}

\bibitem{Choudhury2022}
Abrar Choudhury, Stephen T. Magill, Charlotte D. Eaton, Bria J. H. K. O.
\newblock Meningioma DNA methylation groups identify biological drivers and therapeutic vulnerabilities.
\newblock \textit{Nature Genetics}, 54(5):649--659, 2022.
\newblock \url{https://doi.org/10.1038/s41588-022-01061-8}

\bibitem{Crossref3CA20773}
3CA Catalog Entry: Choudhury et al. 2022
\newblock \textit{Weizmann Institute 3CA}, 2022.
\newblock \url{https://www.weizmann.ac.il/sites/3CA/study-data/20773}

\end{thebibliography}

\end{document}
